\clearpage
\chapter*{第一部分结语}
\addcontentsline{toc}{chapter}{第一部分结语}
\label{part:01-closing}

\begin{center}
  \includegraphics[width=\linewidth,height=.43\textheight,keepaspectratio]{assets/part-art/part-1-closing.png}
  \par\vspace{4pt}
  {\scriptsize\color{Muted}取意于透纳《渡溪》\par}
  \BookEpigraph{part1}
\end{center}

Alice 第一次购买时，账户从 10000 分减到 7000 分，订单和流水也已经保存。她等不到响应，再试一次，账户就变成了 4000 分。沿着这条历史，我们逐步修改了购买服务：先给购买保留身份，再让并发请求争取同一受理位置，最后把原结果与扣款一起保存。客户端依然可能超时，服务也依然可能退出；恢复程序现在有了可以查找和核对的记录。

这套处理从一次确认开始。客户端在发送以前保存完整身份、账户、动作、会话与已确认输入。同一次购买的后续尝试继续引用这条记录，再买一份则另建身份。这个区分关系到用户究竟授权了几次购买。服务收到的身份与输入，承载着调用方已经做出的区分；只有它们正确地对应到确认动作，服务端的受理规则才能对应到实际业务。

身份进入数据库以后，还要取得一项确定的关系：同一完整身份至多对应一条已提交购买。唯一约束检查这项关系，插入结果决定哪个事务继续处理。另一个请求遇到竞争，沿等待、回滚或读取原结果的路径继续。账户的条件更新负责检查余额并扣除金额，订单与流水保存这次成交的归属，完成记录则保存终态和原返回内容。它们各自承担的条件在共同提交处汇合，构成一份完整购买。

原结果保存了这份购买发生时的事实。P1 扣款以后是 7000 分，后来 P2 又把账户减到 4000 分，恢复 P1 仍会读回当次的 7000 分。账户查询负责报告当前余额，订单查询负责报告履约进度。把这些观察分开，调用方才能同时理解先前发生了什么，以及系统现在处于什么状态。结果中的编号、金额与终态也必须和实际订单、流水一致，稳定保存与真实关联共同支撑恢复。

提交和获知结果之间，仍可能相隔很久。一次请求在提交前退出，数据库撤销其内部修改，后来尝试可以重新取得身份；提交已经完成，恢复就沿原身份读取终态。提交确认尚未查明时，恢复者再次经过同一项身份检查，让数据库中的事务结局决定后续分支。客户端保留待确认状态，直到取得原结果或进入明确的核查流程。这段等待期间，已经成立的业务事实继续存在于数据库中。

第一部分的论证依赖明确的范围：购买与本地扣款由同一数据库事务完成，所有相关写入经过规定的仲裁，已确认提交在所讨论的故障后继续保留；恢复使用原身份、原输入和可解释的结果版本，记录在承诺的期限内保持可读。第六章把资金动作限定为购买扣款，因此余额只会保持或下降，余额不足的判断也能在这个模型中安排到终态提交处。加入充值以后，需要重新核对拒绝判断与提交之间的并发关系。操作集合改变，原先依赖的条件也随之改变。

期限同样属于设计的一部分。第一部分在会话有效、记录保留的范围内恢复；期限结束后，第三章的关闭与在途保护边界继续约束清理。持续重试还需要实际机会：数据库恢复可用，旧事务终结，恢复者得到运行，至少一次完整结果能够到达并保存。一次操作已经提交，与客户端最终知道它已经提交，是两个分别需要条件支持的目标。

我们用三种方法检查这些关系。逐步论证说明允许的迁移怎样保持不变量；有界模型展开有限范围内的交错，并让错误方案产生可复核的反例；实际中断实验检查选定位置的进程退出与数据库恢复。它们让设计中的假设、实现中的步骤和观察到的事实能够相互核对。新的故障类型或写入路径进入系统以后，也可以沿这些关系追查：哪一条不变量受影响，哪项前提已经改变，恢复者还能读到什么证据。

接下来，购买服务会把扣款交给另一个系统。对方可以独立完成扣款，本地数据库也有自己的提交结局。第一部分已经连在一起的身份、业务效果与完成记录，将分布在两个提交边界内。我们要继续处理同一个问题：响应丢失以后，恢复者依据哪些持久事实，决定等待、查询或继续执行。
